\section{Canonical Runtime, Docking, and Assembly Semantics}
\label{sec:runtime}

\subsection{Scene and canonical world}

A \code{SceneSpec} is an ordered collection of module placements. It validates unique
instance IDs and known module types; its grid helper creates deterministic names. At
session creation, the backend loads the scene, returns stable handles, and supplies an
initial snapshot. The runtime constructs mutable module, link, joint, connector, and
connection records under \code{WorldState}.

It is useful to qualify the term \emph{event-sourced}. Logical topology and docking
semantics are mutated through events. Backend samples directly update time, poses,
twists, joint state, measured connector frames, and transient lifecycle state. The
event log is therefore a complete record of semantic connection history, not a
lossless physical trajectory recorder.

The world exposes four monotonic revisions: backend sample sequence, topology,
docking-state, and event append count. Derived consumers can invalidate exactly what
they depend upon.

\subsection{Per-step ordering}

\begin{figure}[htbp]
\centering
\begin{tikzpicture}[
  node distance=5mm,
  step/.style={draw,rounded corners,fill=modsimblue!10,draw=modsimblue,
               minimum width=0.78\textwidth,minimum height=8mm,align=center},
  arrow/.style={-{Latex[length=2mm]},thick}
]
  \node[step] (s1) {1. Backend advances physics by $\Delta t$};
  \node[step,below=of s1] (s2) {2. Backend snapshot reports link, joint, connector, and force observations};
  \node[step,below=of s2] (s3) {3. World ingests the fresh sample and resolves connector frames/velocities};
  \node[step,below=of s3] (s4) {4. Overload evaluation requests releases};
  \node[step,below=of s4] (s5) {5. Docking processes releases, candidates, guards, and commits};
  \node[step,below=of s5] (s6) {6. Revisions, events, metrics, and derived views become observable};
  \foreach \a/\b in {s1/s2,s2/s3,s3/s4,s4/s5,s5/s6} \draw[arrow] (\a) -- (\b);
\end{tikzpicture}
\caption{One canonical runtime step. Docking decisions use state reported after the
backend step, and releases are processed before new commits.}
\label{fig:step-pipeline}
\end{figure}

This order prevents decisions against stale geometry. Processing releases before
detection also prevents a connection from breaking and reforming before its release is
visible. A redock cooldown should be authored when two just-released surfaces remain in
capture range.

\subsection{Connector frame and point velocity}

For a connector $C_i$ attached to link $L_i$, an authored numeric pose gives

\begin{equation}
{}^{W}\mathbf T_{C_i} =
{}^{W}\mathbf T_{L_i}\,{}^{L_i}\mathbf T_{C_i}.
\label{eq:connector-transform}
\end{equation}

If the backend materializes and reports the connector frame, that measurement takes
precedence over recomposition. This keeps acceptance consistent with engine geometry,
especially for connectors on articulated child links.

The point velocity includes angular motion:

\begin{equation}
\mathbf v_{C_i} = \mathbf v_{L_i} +
\boldsymbol\omega_{L_i}\times(\mathbf p_{C_i}-\mathbf p_{L_i}).
\label{eq:connector-velocity}
\end{equation}

Without the lever-arm term, a rotating connector could be classified as stationary and
dock above its declared speed limit.

\subsection{Broad phase and deterministic candidates}

Free resolved connectors are indexed in a uniform spatial hash whose cell size is the
largest position tolerance in the pack, with a small positive lower bound. Only nearby
cells produce opportunistic pairs. Explicitly commanded pairs bypass broad-phase
distance rejection so an invalid command produces an informative failure rather than
silence. Pairs are sorted before evaluation.

\subsection{Compatibility}

Type compatibility must be mutual. If type $A$ names $B$ but $B$ does not name $A$,
the pair is rejected as an authoring error rather than interpreted asymmetrically.
Genders pair male--female, hermaphroditic--hermaphroditic, and
genderless--genderless. A connector cannot dock to itself or another connector on the
same module.

\subsection{Geometric and dynamic acceptance}

Acceptance produces a structured result instead of a boolean. For each tolerance $x$,
the effective pair limit is the stricter endpoint declaration,

\begin{equation}
\epsilon_x^{ij}=\min(\epsilon_x^i,\epsilon_x^j).
\label{eq:tolerance}
\end{equation}

Both connector-local capture shapes must contain the relative offset. A sphere uses
Euclidean distance; a box bounds each local coordinate; a cylinder separately bounds
axial and radial displacement.

Two docking axes mate antiparallel, so the axis error is

\begin{equation}
e_a=\cos^{-1}\!\left(\hat{\mathbf a}_i^\top(-\hat{\mathbf a}_j)\right).
\label{eq:axis-error}
\end{equation}

Roll is measured about the shared docking axis using projected connector-frame
reference directions. For a discrete allowed set $\Phi$, the governing error is

\begin{equation}
e_\phi=\min_{\phi_k\in\Phi}
\left|\operatorname{wrap}(\phi-\phi_k)\right|.
\label{eq:roll-error}
\end{equation}

Both endpoints are checked with opposite roll signs. The committed connection stores
the matched orientation index because configuration recognition cannot recover it
reliably from a later noisy pose. Relative connector-point speed is the fourth
criterion. The result reports each measurement, tolerance, margin, and failure reason.

\subsection{Guards and docking policy}

Guards operate after geometry. They reject engaged connectors, cooldown violations,
undefined or conflicting physical intent, unsupported connection types, and missing
commands when the pair is not jointly auto-latching. Closing a loop within an existing
assembly is allowed but warned because rings and lattices are legitimate while
redundant rigid constraints can be accidental.

Alignment policy is either \emph{measured} or \emph{nominal}. Measured alignment freezes
the observed relative pose. Nominal alignment asks the backend to snap connector frames
to their ideal antiparallel/roll pose. If either endpoint requests nominal alignment,
the pair snaps. Physical controller experiments generally prefer measured alignment;
lattice experiments may prefer nominal alignment to avoid accumulated drift.

\subsection{Two-phase physical/logical commit}

\begin{figure}[htbp]
\centering
\begin{tikzpicture}[
  node distance=7mm and 12mm,
  box/.style={draw,rounded corners,align=center,minimum width=35mm,minimum height=9mm},
  decision/.style={diamond,draw,aspect=2.2,align=center,inner sep=1.5pt},
  arrow/.style={-{Latex[length=2mm]},thick}
]
  \node[box,fill=modsimblue!10] (eval) {Recheck acceptance\\and guards};
  \node[box,fill=modsimgold!20,right=of eval] (request) {Create physical\\connection request};
  \node[decision,right=of request] (ok) {Backend\\accepted?};
  \node[box,fill=modsimcyan!15,above right=6mm and 12mm of ok] (commit)
    {Append DockCommitted\\then AssemblyMerged};
  \node[box,fill=modsimorange!12,below right=6mm and 12mm of ok] (fail)
    {Append DockFailed\\with reason/detail};
  \draw[arrow] (eval) -- (request);
  \draw[arrow] (request) -- (ok);
  \draw[arrow] (ok) -- node[above,sloped,font=\scriptsize]{yes} (commit);
  \draw[arrow] (ok) -- node[below,sloped,font=\scriptsize]{no} (fail);
\end{tikzpicture}
\caption{Two-phase connection commit. Canonical topology changes only after the
backend confirms the physical constraint.}
\label{fig:two-phase}
\end{figure}

The backend is first asked to create a constraint. Only a successful returned handle
permits \code{DockCommitted}. If the endpoints previously belonged to separate
components, \code{AssemblyMerged} follows. Backend refusal, unsupported capabilities,
or resource exhaustion records \code{DockFailed}; core never fabricates a connection.

Undocking is symmetric: remove the backend constraint first, then append
\code{UndockCommitted}, then recompute the affected component and append
\code{AssemblySplit} if needed. A permanent connector refuses undocking.

\subsection{Assemblies and topology}

At time $t$, let

\begin{equation}
G_t=(V,E_t), \qquad \mathcal A_t=\operatorname{CC}(G_t),
\label{eq:assemblies}
\end{equation}

where $V$ is the fixed set of module instances and $E_t$ is the set of committed
connections. Parallel edges are preserved even though connected-component membership
does not require multiplicity. Merge joins indexed component sets. Release recomputes
only the affected component rather than the entire world. Assembly identity derives
from membership so it is stable across processes and replay.

\subsection{Events and metrics}

The append-only event vocabulary includes candidate detection, dock success/failure,
undock success/failure, assembly merge/split, and connector overload. Appends receive a
monotonic zero-based sequence. Opportunistic near-miss failures are not logged every
step because they would overwhelm history; explicit failed commands are logged.

Metrics are computed from canonical state and history: module/free/connected counts,
assembly count and maximum size, connector lifecycle counts, active connections,
docking and undocking successes/failures/attempts, merges, splits, overloads, and event
count.

\subsection{Joint commands}

\code{JointCommand} carries a stable module/joint ID, semantic control mode, and scalar
SI target. A batch is validated fully before any backend mutation: IDs must be unique
and known, the pack must declare the mode and applicable limits, the backend must
support the mode, and values must be finite non-booleans. A failure rejects the entire
batch. Commands persist in a backend until replaced or cleared, which makes a complete
batch important for preventing stale effort after a topology change.

\subsection{Overload boundary}

If a backend reports constraint-force magnitude and the two connector policies declare
break-force thresholds, the weaker threshold governs. Exceeding it emits
\code{ConnectorOverloaded} and initiates release. The mock can inject forces for tests;
the current \mujoco{} adapter does not yet report equality force, so physical breakage
is not active there. Normal/shear/bending values stored under connector limits remain
descriptive until a richer load model is implemented.

